\section{The initial shock and all power inequalities}\label{sec:shock}

We examine the internal jump at $r=S$, where the velocity is positive.
For $\lambda\downarrow0$ put
\begin{equation}\label{eq:tangent}
 U_\lambda(\tau,\xi)=u(\lambda\tau,S+\lambda\xi),\quad
 r_\lambda=S+\lambda\xi,\quad
 v_\lambda=v[u](\lambda\tau,r_\lambda),\quad v_0=v[F](S)>0.
\end{equation}
On every fixed half-cylinder, $r_\lambda^2\to S^2$ and
$v_\lambda\to v_0$ uniformly, and $A/2\leq U_\lambda\leq U$ for
small $\lambda$. The initial data tend in every finite local $L^q$ to
the step $d_*(\xi)=1$ for $\xi<0$ and $d_*(\xi)=U$ for $\xi>0$.

\subsection{Compactness and identification of the tangent}

The scaled mass flux is $r_\lambda^2v_\lambda f(U_\lambda)$ and the
scaled entropy flux is $r_\lambda^2v_\lambda Q_\eta(U_\lambda)$.
For the linear fractional term, let
\[
 \widetilde g(t,r)=rF_\eta(u(t,|r|))\quad(r\in\R),\qquad
 g_\lambda(\tau,\xi)=\widetilde g(\lambda\tau,S+\lambda\xi).
\]
It is the pullback of the odd physical radial profile. Its supremum norm
is bounded uniformly, although its support in $\xi$ grows. The scaled term is
\begin{equation}\label{eq:scaled-source}
 \lambda^{1-\alpha}r_\lambda L_{1,\alpha}g_\lambda.
\end{equation}
Transferring $L_{1,\alpha}$ to a compact smooth test bounds its pairing by
$\lambda^{1-\alpha}\|g_\lambda\|_\infty
\|L_{1,\alpha}(r_\lambda\varphi)\|_1\to0$.
The last norm is uniformly bounded by \eqref{eq:fractional-L1}.

For a convex entropy the other scaled sources are positive measures.
Testing the entropy balance with a nonnegative cutoff equal to one on a
fixed rectangle bounds their mass there: the entropy, flux and initial
terms are bounded, and \eqref{eq:scaled-source} has a bounded pairing.
In particular the quadratic measure has a uniform local mass bound.
The exact coefficients for Lemma~\ref{lem:increments} at zero viscosity are
\[
 \omega_\lambda=r_\lambda^2,\qquad
 b_\lambda=\lambda^{1-\alpha}r_\lambda,\qquad
 g=g_\lambda,\qquad\eps=0.
\]
The weights are uniformly positive and Lipschitz on each fixed rectangle.
The spatial H\"older seminorm of $v_\lambda$ gains $\lambda^\beta$ from
\eqref{eq:velocity-holder}; no time derivative of $v_\lambda$ is required.
Thus, on a smaller fixed rectangle,
\begin{equation}\label{eq:scaled-increments}
 \int|U_\lambda(\tau,\xi+h)-U_\lambda(\tau,\xi)|^4\,\dd\xi\dd\tau
 \leq K|h|^{1-\alpha}.
\end{equation}
This is a fourth-power bound; the corresponding $L^4$ norm has exponent
$(1-\alpha)/4$. It uses the uniform supremum of $g_\lambda$, without
a bound on its growing spatial $L^1$ norm.

Cauchy--Schwarz turns \eqref{eq:scaled-increments} into a squared
$L^2$ increment bound $K|h|^{(1-\alpha)/2}$. After spatial localization,
the translation formula for Sobolev norms gives a uniform
$L^2_\tau H^\sigma_\xi$ bound for each
$0<\sigma<(1-\alpha)/4$. The weighted mass equation bounds
$\partial_\tau(r_\lambda^2U_\lambda)$ in $L^1_\tau H^{-1}_\xi$.
Use \cite[Lemma~7.5]{Stan2016} on bounded intervals with
$X=H^\sigma$, $B=L^2$, $Y=H^{-1}$ and exponent two, then divide by the
smooth positive weight. A diagonal extraction gives, from every scale
sequence, a tangent $Z$ with strong local $L^2$ convergence.
If compactness is first applied away from $\tau=0$, the cap discards an
initial strip with squared $L^2$ mass $O(\delta)$ uniformly, so the
convergence extends to compact half-cylinders. Bounded states upgrade it
to every finite local $L^q$.

Passing the scaled mass and convex entropy balances yields
\begin{equation}\label{eq:scalar-limit}
 Z_\tau+v_0\partial_\xi f(Z)=0,\qquad
 \partial_\tau\eta(Z)+v_0\partial_\xi Q_\eta(Z)\leq0,
\end{equation}
with the actual step initial mass and entropy terms. Subtracting affine
functions extends the normalized inequalities to all smooth convex
entropies. The mass equation selects compact-test continuous pairings.
The localized quadratic entropy inequality gives
\[
 \limsup_{\tau\downarrow0}\int\chi Z(\tau)^2
 \leq\int\chi d_*^2
\]
for nonnegative compact $\chi$, first at good times and then by weak
lower semicontinuity for the representative. Its weak initial trace and
bounded states give $\int\chi Z(\tau)d_*\to\int\chi d_*^2$, by smooth
approximation of $\chi d_*$. Expanding the square proves the strong
local $L^2$, hence local $L^1$, initial trace. This supplies the scalar
initial trace without assuming a strong trace at the unscaled inner jump.

The bounded scalar uniqueness theorem
\cite[Section~6.1, Theorem~4]{Bressan2009} now identifies the tangent.
To match its smooth-flux hypothesis, extend $v_0k^a$ smoothly to $\R$
outside a neighborhood of $[A/2,U]$, using a cutoff supported in the
positive half-line. All compared states stay in that interval, so the
equation and entropy flux derivatives are unchanged. The step itself
need not be integrable; the theorem treats bounded data, and the compared
initial difference is zero. The solution is
\begin{equation}\label{eq:shock}
 Z_*(\tau,\xi)=\begin{cases}1,&\xi<\sigma_0\tau,\\
 U,&\xi>\sigma_0\tau,\end{cases}
 \qquad \sigma_0=v_0\frac{U^a-1}{U-1}=\frac{v_0}{11}.
\end{equation}
It satisfies Rankine--Hugoniot. Recall the chord deficit
$H_U(k)=f(k)-f(1)-[(f(U)-f(1))/(U-1)](k-1)$.
For every convex $C^2$ entropy, integration by parts in $k$, using
$H_U(1)=H_U(U)=0$, gives its exact
entropy distribution
\begin{equation}\label{eq:shock-production}
 \partial_\tau\eta(Z_*)+v_0\partial_\xi Q_\eta(Z_*)
 =-v_0\left(\int_1^U\eta''(k)H_U(k)\,\dd k\right)
             \delta_{\xi=\sigma_0\tau}\,\dd\tau.
\end{equation}
Its sign is nonpositive. Uniqueness identifies every subsequential tangent
with \eqref{eq:shock}. Since every scale sequence is relatively compact,
the entire family $U_\lambda$ converges to $Z_*$ in every finite local
$L^q$.

\subsection{Dissipation at higher densities}

Fix any $A<B<1$. Choose $0\leq\zeta\leq1$ smooth, zero for $k\leq B$
and one on $[1,U]$, and define a normalized convex entropy by
$\eta''(k)=\zeta(k)/k$, with zero extension near zero. It has bounded
second derivative on $[0,U]$, so \eqref{eq:limit-local-entropy} applies.
Its shock coefficient in \eqref{eq:shock-production} is $-v_0J$, where
$J=\int_1^U H_U(k)/k\,\dd k$.

Under $t=\lambda\tau$, $r=S+\lambda\xi$, the radial viscosity measure is
\[
 \bar\nu_\lambda=\lambda^{-1}
 (t,r,k\mapsto t/\lambda,(r-S)/\lambda,k)_\#\bar\nu.
\]
The factor $\lambda^{-1}$ follows from the first derivative scaling
and the two-dimensional time-radius Jacobian. In the exact entropy
balance, the linear fractional term vanishes by
\eqref{eq:scaled-source}, while \eqref{eq:tube-defect} makes the separated
nonlocal defect $o(\lambda)$ on each physical tube of width and length
$O(\lambda)$. Full tangent convergence therefore gives, weakly against
compact tests with $\tau>0$,
\begin{equation}\label{eq:viscosity-shock}
 \int\frac{\zeta(k)}k\,\bar\nu_\lambda(\dd\tau,\dd\xi,\dd k)
 \ \rightharpoonup\ S^2v_0J\,
              \delta_{\xi=\sigma_0\tau}\,\dd\tau.
\end{equation}
The local positive measure bounds above justify this weak interpretation.
In particular the nonlocal defect has not been counted as viscosity
production.

Test \eqref{eq:viscosity-shock} by a nonnegative cutoff at most one,
supported in $0<\tau<1$ and a fixed spatial interval containing the
shock segment, and equal to one on the shock for
$\delta\leq\tau\leq1-\delta$. Since
$\zeta(k)/k\leq\mathbf1_{[B,U]}(k)/k$, positivity and
\eqref{eq:radial-nu} imply, on sending $\delta\downarrow0$,
\begin{equation}\label{eq:high-density}
 \liminf_{t\downarrow0}\frac1t
 \int_{(0,t)\times\R^3\times[B,U]}\frac1k\,\dd\nu
 \geq4\pi S^2v_0J=W[F;S]J.
\end{equation}
This is a full liminf, since the tangent convergence used no favorable
scale sequence.

The outer weight $W[u(t);R(t)]$ is positive and continuous at zero, with
value $W[F;R_0]$. Recall that $G_A=\int_0^A g_A(k)/k\,\dd k$ and
the fixed strict margin is $W[F;R_0]G_A<W[F;S]J$. This margin and
\eqref{eq:high-density} therefore supply one $0<T_1<T$ such that
\begin{equation}\label{eq:payment}
 \int_{(0,t)\times\R^3\times[B,U]}\frac1k\,\dd\nu
 \geq G_A\int_0^tW[u(r);R(r)]\,\dd r,
 \qquad0<t\leq T_1.
\end{equation}
For example choose an error smaller than one third of the positive gap in
\eqref{eq:strict-margin}, apply the definition of liminf on one interval
and continuity of the averaged outer weight on another, and intersect
them. This time is chosen independently of every later power exponent.

\subsection{The exact power coefficient and every terminal time}

For $b>0$ and $r,z\geq0$, Cauchy--Schwarz on the interval between the
two states gives
\begin{equation}\label{eq:scalar-power}
 (r^b-z^b)(r-z)\geq\frac{4b}{(b+1)^2}
       (r^{(b+1)/2}-z^{(b+1)/2})^2.
\end{equation}
Indeed apply it to $\int_z^r k^{(b-1)/2}\dd k$, and then use a positive
lower endpoint and monotone convergence if one state is zero.
Set $b=a+q-1$, so
$\Psi_{k^q}(k)=q(q-1)k^b/b$ by \eqref{eq:psi-power}.
Combining this formula, \eqref{eq:scalar-power} and the positive pair form gives
\begin{equation}\label{eq:power-fractional}
 \mathcal B_\alpha(\Psi_{k^q},w)
 \geq\frac{4q(q-1)}{(a+q)^2}
 \|(-\Delta)^{\alpha/4}w^{(a+q)/2}\|_2^2.
\end{equation}
This holds as a comparison of nonnegative extended integrals. For
$q=2-a$ the scalar inequality is equality; no strict fractional remainder
is available at that exponent.

Take, for example, $B=(A+1)/2$. On $k\geq B$,
$k^{q-2}\geq B^{q-1}/k$. The viscosity moment in
\eqref{eq:limit-power-cost} is therefore at least $B^{q-1}$ times the
left side of \eqref{eq:payment}. The positive boundary contribution in that balance is at most
$A^{q-1}G_A\int_0^tW[u;R]$ by \eqref{eq:density-cost}.
Both terms carry the same factor $q(q-1)$ in \eqref{eq:limit-power-cost}.
Since $B>A$, the same measure and the same $T_1$ make the viscosity
contribution dominate the boundary contribution for every $q>1$.
At good times the result, followed by
\eqref{eq:power-fractional}, is
\begin{equation}\label{eq:short-power}
 \|u(t)\|_q^q+\frac{4q(q-1)}{(a+q)^2}
 \int_0^t\|(-\Delta)^{\alpha/4}u^{(a+q)/2}\|_2^2\,\dd r
 \leq\|F\|_q^q.
\end{equation}
For any terminal $t\leq T_1$, approach it by increasing good times.
Bounded-test weak $L^1$ continuity gives lower semicontinuity of the
convex power norm, and the nonnegative time integral increases to its
full value. This extends \eqref{eq:short-power} to every terminal time.
The dense-exponent passage described after \eqref{eq:limit-power-cost}
gives every real $q>1$ without changing the measure or time interval.
Together with \eqref{eq:limit-pressure}, the distributional equation,
mass, cap and \eqref{eq:weak-time}, this proves all required energy weak
solution properties on $[0,T_1]$. Its support is contained in the finite
$C^1$ radius $R(t)$ there.
